% Pitch Predicta — 5 minutos · Hackathon
% Compilar: pdflatex predicta_pitch.tex (2x)
\documentclass[aspectratio=169,11pt]{beamer}
\usepackage[utf8]{inputenc}
\usepackage[T1]{fontenc}
\usepackage[brazil]{babel}
\usepackage{tikz}
\usetikzlibrary{arrows.meta,positioning,shapes.geometric,calc}
\usepackage{booktabs}

% ---------- Tema ----------
\usetheme{default}
\usecolortheme{seagull}
\definecolor{predicta}{HTML}{4F46E5}
\definecolor{verde}{HTML}{059669}
\definecolor{laranja}{HTML}{DC6B18}
\definecolor{grafite}{HTML}{1F2937}
\definecolor{cinza}{HTML}{64748B}
\setbeamercolor{frametitle}{fg=predicta}
\setbeamercolor{title}{fg=predicta}
\setbeamercolor{structure}{fg=predicta}
\setbeamertemplate{navigation symbols}{}
\setbeamertemplate{footline}{%
  \hfill\scriptsize\textcolor{cinza}{Predicta · Hackathon 2026 · \insertframenumber/\inserttotalframenumber}\hspace{2mm}\vspace{1mm}}
\setbeamerfont{frametitle}{series=\bfseries,size=\Large}

\title{\Huge Predicta}
\subtitle{Tarifa dinâmica inteligente: prever a demanda, achatar o pico}
\author{Equipe Predicta}
\date{}

\begin{document}

% ---------- Slide 0 · Capa ----------
\begin{frame}
  \titlepage
  \centering
  \textcolor{cinza}{Demanda energética, clima e operação do sistema}
\end{frame}

% ---------- Slide 1 · Problema (0:00–0:40) ----------
\begin{frame}{O sistema é dimensionado pelo pior dia}
  \begin{columns}[T]
    \begin{column}{0.55\textwidth}
      \vspace{2mm}
      \begin{itemize}\setlength\itemsep{3mm}
        \item \textbf{Demanda $\neq$ consumo}: dois dias com o mesmo total de kWh podem exigir condições operativas muito diferentes.
        \item \textbf{Calor extremo concentra carga} em poucas horas — o pico define o custo de rede, geração e reserva.
        \item \textbf{O consumidor não tem incentivo} para sair do pico: a tarifa que ele vê é plana.
      \end{itemize}
      \vspace{3mm}
      \begin{beamercolorbox}[rounded=true,sep=3mm]{block body}
        \centering\textbf{\textcolor{laranja}{``Dimensionado pelo pior dia,\\ cobrado como se todo dia fosse igual.''}}
      \end{beamercolorbox}
    \end{column}
    \begin{column}{0.42\textwidth}
      % Espaço para gráfico: curva de carga com pico destacado (screenshot /territorio ou ilustração)
      \centering
      \begin{tikzpicture}[scale=0.9]
        \draw[->,thick,cinza] (0,0) -- (5.2,0) node[below left,font=\scriptsize]{hora};
        \draw[->,thick,cinza] (0,0) -- (0,3.2) node[left,font=\scriptsize]{MW};
        \draw[very thick,predicta] plot[smooth] coordinates {(0.2,1.2)(1,0.9)(2,1.3)(3,2.0)(3.8,2.8)(4.4,2.2)(5,1.4)};
        \draw[dashed,laranja,thick] (0,2.8) -- (5.2,2.8) node[above left,font=\scriptsize\bfseries,text=laranja]{pico dimensiona tudo};
      \end{tikzpicture}
    \end{column}
  \end{columns}
\end{frame}

% ---------- Slide 2 · Solução: os dois motores (0:40–1:50) ----------
\begin{frame}{A solução: dois motores, um sinal auditável}
  \centering
  \begin{tikzpicture}[
      font=\scriptsize,
      fonte/.style={draw=cinza,rounded corners=1pt,fill=white,minimum width=2.55cm,minimum height=0.52cm,align=center},
      motor/.style={draw=predicta,very thick,rounded corners=3pt,fill=predicta!8,minimum width=3.1cm,minimum height=2.9cm,align=center},
      saida/.style={draw=verde,thick,rounded corners=2pt,fill=verde!8,minimum width=2.9cm,minimum height=1.15cm,align=center},
      seta/.style={-{Stealth[length=2.5mm]},thick,grafite}]
    % Fontes de dados (colocar logos aqui)
    \node[fonte] (f1) at (0,2.4) {ONS · Carga};
    \node[fonte] (f2) at (0,1.7) {ONS · Geração};
    \node[fonte] (f3) at (0,1.0) {ONS · DESSEM};
    \node[fonte] (f4) at (0,0.3) {ANEEL · Tarifas};
    \node[fonte] (f5) at (0,-0.4) {ERA5-Land · Clima 0,1$^{\circ}$};
    \node[fonte] (f6) at (0,-1.1) {Ativos geoespaciais};
    \node[fonte,dashed] (f7) at (0,-1.8) {Open-Meteo forecast$^{*}$};
    % Motor 1
    \node[motor] (m1) at (4.0,0.3) {\textbf{\normalsize Motor 1 · SIN}\\[1.5mm] Previsão de demanda\\ 24h por subsistema\\ (p10 · p50 · p90)\\[1mm] pressão de oferta\\ exposição climática};
    % Sinal
    \node[draw=laranja,very thick,rounded corners=2pt,fill=laranja!10,minimum height=0.8cm,align=center] (sig) at (7.6,0.3) {\textbf{system\_signal\_v1}\\ sinal horário auditável};
    % Motor 2
    \node[motor] (m2) at (11.2,0.3) {\textbf{\normalsize Motor 2 · Tarifa}\\[1.5mm] \textbf{Guardrails:}\\ piso/teto 0,85–1,30$\times$\\ rampa máx. $\pm$10\%/h\\ fatura média neutra\\ trava de +20\%/fatura\\ falhou o dado? $\to$ 1,0$\times$};
    % Saídas
    \node[saida] (s1) at (14.6,1.3) {\textbf{Consumidor}\\ quando gastar menos};
    \node[saida] (s2) at (14.6,-0.7) {\textbf{Distribuidora}\\ curva achatada};
    % Setas
    \foreach \f in {f1,f2,f3,f4,f5,f6,f7} \draw[seta,cinza] (\f.east) -- ($(m1.west)+(0,0)$);
    \draw[seta,very thick,laranja] (m1) -- (sig);
    \draw[seta,very thick,laranja] (sig) -- (m2);
    \draw[seta,verde] (m2.east) -- (s1.west);
    \draw[seta,verde] (m2.east) -- (s2.west);
  \end{tikzpicture}

  \vspace{1mm}
  {\scriptsize\textcolor{cinza}{$^{*}$previsão meteorológica operacional: em integração (roadmap). Todas as demais fontes já alimentam o MVP.}}
\end{frame}

% ---------- Slide 3 · Demo (1:50–2:30) ----------
\begin{frame}{Demo ao vivo · o produto funcionando}
  \begin{columns}[T]
    \begin{column}{0.55\textwidth}
      \vspace{4mm}
      \textbf{Roteiro (40s):}
      \begin{enumerate}\setlength\itemsep{2mm}
        \item Escolher a distribuidora no mapa real de concessões (ANEEL);
        \item Simular um dia real: \textbf{DESSEM $\times$ Predicta $\times$ Realizado};
        \item Card do consumidor: \textbf{economia potencial R\$/mês};
        \item Aba da carteira: \textbf{alívio de pico em MW} com a adesão simulada.
      \end{enumerate}
    \end{column}
    \begin{column}{0.42\textwidth}
      \vspace{4mm}
      % Substituir pelo screenshot de backup da tela /produto/
      \begin{beamercolorbox}[rounded=true,sep=8mm]{block body}
        \centering\textcolor{cinza}{[screenshot de backup\\ da tela /produto/]}
      \end{beamercolorbox}
    \end{column}
  \end{columns}
\end{frame}

% ---------- Slide 4 · Confiança (2:30–3:10) ----------
\begin{frame}{Dá para confiar? Cada número tem origem}
  \begin{columns}[T]
    \begin{column}{0.31\textwidth}
      \centering
      \tikz\node[circle,draw=predicta,very thick,minimum size=9mm,font=\Large\bfseries,text=predicta]{1};\\[2mm]
      \textbf{Fontes oficiais}\\[1mm]
      {\small Tudo vem de ONS, ANEEL e bases climáticas públicas. Nada inventado.}
    \end{column}
    \begin{column}{0.31\textwidth}
      \centering
      \tikz\node[circle,draw=predicta,very thick,minimum size=9mm,font=\Large\bfseries,text=predicta]{2};\\[2mm]
      \textbf{``Nota fiscal'' de dados}\\[1mm]
      {\small Cada tarifa guarda o carimbo de origem de cada insumo. Qualquer valor pode ser reproduzido e auditado depois.}
    \end{column}
    \begin{column}{0.31\textwidth}
      \centering
      \tikz\node[circle,draw=predicta,very thick,minimum size=9mm,font=\Large\bfseries,text=predicta]{3};\\[2mm]
      \textbf{Testado contra a realidade}\\[1mm]
      {\small Modelo congelado e confrontado com 30 dias reais que ele nunca viu.}
    \end{column}
  \end{columns}
  \vspace{6mm}
  \centering
  \begin{beamercolorbox}[rounded=true,sep=4mm]{block body}
    \centering
    {\LARGE\textbf{\textcolor{verde}{Erro de previsão: metade do método tradicional}}}\\[2mm]
    {\normalsize 2,3\% vs 4,5\% no Sudeste/Centro-Oeste · 3,3\% vs 6,7\% no Sul}\\[1mm]
    {\footnotesize\textcolor{cinza}{e nas horas de calor extremo — quando mais importa — o erro se mantém em $\approx$2\%}}
  \end{beamercolorbox}
\end{frame}

% ---------- Slide 5 · Modelo de negócio (3:10–3:50) ----------
\begin{frame}{Quem usa, quem paga, quem ganha}
  \centering
  \begin{tikzpicture}[font=\small,
      caixa/.style={draw=predicta,thick,rounded corners=3pt,fill=predicta!6,minimum width=4.2cm,minimum height=2.5cm,align=center},
      seta/.style={-{Stealth[length=2.5mm]},thick,cinza}]
    \node[caixa] (dist) at (0,0) {\textbf{Distribuidora}\\[1.5mm] paga \textbf{setup} + \textbf{fee de serviço}\\[1mm] ganha achatando a curva:\\ mais venda na ociosidade,\\ menos estresse no pico};
    \node[caixa] (cons) at (6.4,0) {\textbf{Consumidor}\\[1.5mm] previsão 24h \textbf{grátis}\\ via distribuidora\\[1mm] horizonte estendido +\\ otimização = \textbf{fee premium} opcional};
    \node[caixa,fill=verde!8,draw=verde] (sys) at (12.8,0) {\textbf{Sistema elétrico}\\[1.5mm] pico menor $\to$ menos\\ despacho térmico caro,\\ menos investimento\\ em rede ociosa};
    \draw[seta] (dist) -- (cons);
    \draw[seta] (cons) -- (sys);
  \end{tikzpicture}
  \vspace{4mm}

  {\small\textcolor{cinza}{Receita recorrente para a Predicta · payback da distribuidora vem do custo evitado de pico.}}
\end{frame}

% ---------- Slide 6 · Viabilidade e roadmap (3:50–4:25) ----------
\begin{frame}{MVP funcional hoje, piloto amanhã}
  \begin{columns}[T]
    \begin{column}{0.48\textwidth}
      \textbf{\textcolor{verde}{O que já roda hoje}}
      \begin{itemize}\setlength\itemsep{1.5mm}
        \item Pipeline completo: dados oficiais $\to$ previsão $\to$ tarifa $\to$ web;
        \item 4 subsistemas do SIN, histórico 2023--2026;
        \item Simulador de cliente e de carteira com dados reais;
        \item API REST pronta para integrar ao billing.
      \end{itemize}
    \end{column}
    \begin{column}{0.48\textwidth}
      \textbf{\textcolor{laranja}{Próximos passos (honestos)}}
      \begin{itemize}\setlength\itemsep{1.5mm}
        \item Previsão meteorológica ao vivo (código pronto, falta operação);
        \item Operação agendada 24/7;
        \item \textbf{Piloto com 1 distribuidora} em ambiente regulatório experimental (sandbox ANEEL).
      \end{itemize}
    \end{column}
  \end{columns}
  \vspace{4mm}
  \centering
  {\small\textcolor{cinza}{Fase 1: piloto shadow (tarifa simulada em paralelo à real) $\to$ Fase 2: opt-in com clientes voluntários $\to$ Fase 3: escala.}}
\end{frame}

% ---------- Slide 7 · Inovação (4:25–4:50) ----------
\begin{frame}{Por que isso é diferente}
  \begin{table}
    \centering\small
    \begin{tabular}{lccc}
      \toprule
       & \textbf{Bandeiras} & \textbf{Tarifa branca} & \textbf{Predicta} \\
      \midrule
      Granularidade & mensal & 3 postos fixos & \textbf{horária} \\
      Antecipação & reativa & estática & \textbf{preditiva 24h} \\
      Considera clima & não & não & \textbf{sim} \\
      Proteção ao consumidor & — & — & \textbf{guardrails} \\
      Localizada por região & não & não & \textbf{sim} \\
      \bottomrule
    \end{tabular}
  \end{table}
  \vspace{3mm}
  \centering
  {\small Dois motores desacoplados por contrato de dados: a distribuidora pluga o Motor 2\\ no billing existente sem reconstruir nada.}
\end{frame}

% ---------- Slide 8 · Fechamento (4:50–5:00) ----------
\begin{frame}{}
  \centering
  \vspace{8mm}
  {\Huge\textbf{\textcolor{predicta}{Predicta}}}\\[4mm]
  {\Large A demanda vai continuar mudando com o clima.\\[1mm] A tarifa precisa aprender a acompanhar.}\\[8mm]
  {\large\textcolor{verde}{\textbf{Buscamos uma distribuidora parceira para o piloto.}}}\\[6mm]
  {\small\textcolor{cinza}{Equipe Predicta · contato}}
\end{frame}

\end{document}
% Nota: os círculos numerados do slide 4 são marcadores — substituir por ícones (ex: pacote fontawesome5:
% \usepackage{fontawesome5} e \faDatabase, \faStamp, \faCheckDouble) ou imagens da identidade visual.
